\documentclass[french,11pt]{book}
\usepackage{teach}
\usepackage[french]{babel}
\usepackage{amsmath,amssymb}
\usepackage{array}
\newcommand{\R}{\mathbb{R}}
\begin{document}
\begin{center}
{\Large \textbf{Devoir surveillé -- Fonction inverse}}\\[2mm]
Terminale technologique
\end{center}

\section*{Devoir surveillé}
\begin{enumerate}
  \item Soit $f(x)=\frac{1}{x}$. Calculer $f(-5)$, $f(0,25)$ et $f\left(\frac{2}{3}\right)$.
  \item Résoudre $\frac{1}{x}=2,5$ puis interpréter graphiquement la solution.
  \item On étudie le coût unitaire $C(x)=\frac{360}{x}$ pour $x$ objets produits. Calculer $C(30)$ et déterminer à partir de combien d'objets le coût unitaire devient inférieur ou égal à 9 euros.
\end{enumerate}

\end{document}
